% © 2025 Ing. David Jaroš — CC BY-NC-ND 4.0
%
% This work is licensed under a Creative Commons Attribution-NonCommercial-NoDerivatives 
% 4.0 International License (CC BY-NC-ND 4.0).
%
% License History: Earlier drafts (up to v0.3) were released under CC BY 4.0. 
% From v0.4 onward, all material is released under CC BY-NC-ND 4.0 to protect 
% the integrity of the theoretical work during ongoing academic development.
%
% See LICENSE.md for full license text.

\ifdefined\INCLUDEMODE
\else
  \documentclass[12pt]{article}
  \usepackage[a4paper, margin=2.5cm]{geometry}
  \usepackage{amsmath, amssymb, amsthm}
  \usepackage{hyperref}
  \usepackage{xcolor}
  \usepackage{mdframed}
  \usepackage{booktabs}

  \newmdenv[backgroundcolor=green!15, linecolor=green!60!black, linewidth=1.5pt]{resultbox}
  \newmdenv[backgroundcolor=blue!10, linecolor=blue!60, linewidth=1.5pt]{predbox}
  \newmdenv[backgroundcolor=yellow!20, linecolor=orange, linewidth=1.5pt]{gapbox}

  \title{\textbf{Step 6: QED Reproducibility Summary at $\varphi = \mathrm{const}$}}
  \author{David Jaroš}
  \date{2026-03-06}

  \begin{document}
  \maketitle
\fi

\section{Step 6: QED Reproducibility Summary at $\varphi = \mathrm{const}$}
\label{sec:qed_summary}

\textbf{Task reference:} UBT\_v29\_task7, step6\_qed\_summary.
\textbf{Date:} 2026-03-06.

\subsection{Status Table}

\begin{center}
\renewcommand{\arraystretch}{1.4}
\begin{tabular}{lllp{5.5cm}}
\toprule
\textbf{QED Observable} & \textbf{$\varphi=0$ status} & \textbf{$\varphi=\mathrm{const}$ status} & \textbf{Notes} \\
\midrule
Photon massless & \checkmark & \textbf{PROVED} (Step~1) & $q_\varphi=0$; no Higgs for EM \\
Electron mass $m_e$ & \checkmark & \textbf{SKETCH} (Step~2) & $m_e=yv$; gaps Y1, Y2 remain \\
$\alpha(\mu)$ running & \checkmark & \textbf{PROVED} (Step~3) & $\delta B(\varphi)=0$ at 1-loop \\
$g$-factor Schwinger $a_e=\alpha/(2\pi)$ & \checkmark & \textbf{PROVED} (Step~4) & Sub-leading $\varphi$-corrections $\sim 10^{-16}$ \\
Lamb shift $1057.8\,\mathrm{MHz}$ & \checkmark & \textbf{SKETCH} (Step~5) & Structural; Gap~L1 open \\
\bottomrule
\end{tabular}
\end{center}

\subsection{Overall Assessment}

\begin{resultbox}
\textbf{Minimum success criterion: MET.}
$U(1)_\mathrm{EM}$ is proved unbroken at $\varphi = \mathrm{const}$ (Step~1).
Two of five observables are fully proved (photon mass zero; Schwinger term);
one ($\alpha(\mu)$ running) is proved at one loop.
Two observables (electron mass, Lamb shift) are at SKETCH level with documented gaps.
\end{resultbox}

\textbf{QED reproducibility at $\varphi = \mathrm{const}$: SUBSTANTIALLY PROVED at one loop.}

The key structural result is:
\begin{enumerate}
  \item $\varphi$ is a $U(1)_\mathrm{EM}$-neutral scalar (Step~1, proved).
  \item Because $q_\varphi = 0$, the $\varphi$-background does not contribute any new
        one-loop diagram to the photon self-energy, vertex correction, or electron
        self-energy.
  \item All QED one-loop quantities in the $\varphi = \mathrm{const}$ background are
        identical to the $\varphi = 0$ quantities up to the replacement
        $m_e \to yv$ in finite parts, which does not affect UV-divergence-controlled
        observables ($\beta$-function, Schwinger term) at one loop.
\end{enumerate}

\subsection{Open Gaps}

\begin{center}
\begin{tabular}{lll}
\toprule
\textbf{Gap} & \textbf{Description} & \textbf{Priority} \\
\midrule
Y1 & Derive Yukawa $y$ from $S[\Theta]$ & HIGH \\
Y2 & Derive VEV $v$ from $V_\mathrm{eff}(\theta_0)$ & HIGH \\
L1 & Explicit UBT Lamb-shift computation (path-integral) & MEDIUM \\
\bottomrule
\end{tabular}
\end{center}

\subsection{UBT Predictions from $\varphi$-Background}

\begin{predbox}
\textbf{Prediction P-QED-1 (two-loop; sub-leading):}
At two-loop order, the $\varphi$-background generates a tiny correction to the
fine structure constant running:
\[
  \frac{\delta\alpha}{\alpha}\bigg|_\varphi
  \;\sim\; \frac{y_e^2}{16\pi^2}\cdot\frac{\alpha}{\pi}
  \;\approx\; 10^{-21}.
\]
This is unobservably small with current technology but in principle constitutes a
testable prediction of UBT in regions of different cosmological $\varphi$ background.
\end{predbox}

\begin{predbox}
\textbf{Prediction P-QED-2 (two-loop; sub-leading):}
The electron anomalous magnetic moment receives a correction
\[
  \delta a_e^{(\varphi)} \;\sim\; \frac{y_e^2}{16\pi^2}\cdot\frac{\alpha}{2\pi}
  \;\approx\; 10^{-16},
\]
which is below current experimental sensitivity ($\sim 10^{-13}$).
\end{predbox}

\subsection{Derivation Index Entries Added}

The following results from this task chain are to be added to \texttt{DERIVATION\_INDEX.md}:

\begin{itemize}
  \item $U(1)_\mathrm{EM}$ unbroken at $\varphi = \mathrm{const}$: \textbf{Proven [L0]},
        \texttt{consolidation\_project/qed\_phi\_const/step1\_u1\_protection.tex}
  \item $\alpha(\mu)$ running at $\varphi = \mathrm{const}$: \textbf{Proven [L1]},
        $\delta B=0$ (one-loop); \texttt{step3\_beta\_function.tex}
  \item Schwinger $a_e = \alpha/(2\pi)$ at $\varphi = \mathrm{const}$: \textbf{Proven [L1]},
        \texttt{step4\_schwinger\_term.tex}
  \item Electron mass $m_e = yv$ at $\varphi = \mathrm{const}$: \textbf{Sketch [L1]},
        \texttt{step2\_electron\_mass.tex} (Gaps Y1, Y2)
  \item Lamb shift at $\varphi = \mathrm{const}$: \textbf{Sketch [L1]},
        \texttt{step5\_lamb\_shift.tex} (Gap L1)
\end{itemize}

\subsection{Relation to Existing QED Limit}

Previously, UBT validated QED only at $\varphi \to 0$ (vacuum limit):
$B_0^\mathrm{QED} = 2\pi/3$ recovered (see \texttt{DERIVATION\_INDEX.md §Fine Structure Constant}).
The present task extends this validation to the physically realistic case
$\varphi = \mathrm{const} \neq 0$.

The extension does not modify any existing results; it strengthens them by proving that
the QED sector of UBT is robust against a constant scalar background.

\subsection{Numerical Verification}

A Python script \texttt{tools/verify\_qed\_phi\_const.py} accompanies this derivation
chain.  It verifies:
\begin{itemize}
  \item That $\delta B(\varphi) = 0$ for constant $\varphi$ (analytic check);
  \item That $a_e(\varphi) = \alpha/(2\pi)$ numerically (to $10^{-15}$ relative precision);
  \item That the two-loop $\varphi$-corrections are bounded as predicted.
\end{itemize}

Run with: \texttt{python tools/verify\_qed\_phi\_const.py}

\ifdefined\INCLUDEMODE\else
  \end{document}
\fi
